% ------------------------------------------------------------------
%  Clustering method: K--Means.
%  Written from template/method_template.tex; the part order is fixed.
%  Code counterpart: xxcluster/cluster/partitional/sse_based/kmeans.py
% ------------------------------------------------------------------
\subsubsection{K–Means}
\label{sec:tech:kmeans}

\paragraph{Overview}
\label{sec:tech:kmeans:overview}
K–Means is the canonical crisp partitional method, belonging to the SSE–optimised family of \textbf{Sect. \ref{sec:background:families}}.
It represents every cluster by a centroid and assigns each observation to the nearest one, so a partition is described by $|\mathcal C|$ prototypes rather than by the observations \cite{ref_10}. It serves as the baseline the other methods are compared against \cite{ref_1, ref_5}, and has precedent on operational water data \cite{ref_13, ref_17}.

\paragraph{Assumptions and applicability}
\label{sec:tech:kmeans:assumptions}
Assignment to the nearest centroid induces a Voronoi partition of the feature space, hence clusters are assumed convex, roughly spherical, and of comparable size and density \cite{ref_10, ref_11}.
It further assumes $|\mathcal C|$ known in advance, no noise label, and features on comparable scales, the objective below not being scale–invariant, where these fail, the partition remains well defined but ceases to be meaningful.

\paragraph{Formulation}
\label{sec:tech:kmeans:formulation}
For fixed $|\mathcal C|$, K–Means minimises the within–cluster sum of squared error
\begin{align}
    \min_{\mathcal C} ~ \sum_{i = 1}^{|\mathcal C|} ~ \sum_{\mathbf x \in \mathcal C_i} d(\mathbf x, \boldsymbol \mu_i)^2,
    \qquad \text{where} \qquad
    \boldsymbol \mu_i = \frac{1}{|\mathcal C_i|} \sum_{\mathbf x \in \mathcal C_i} \mathbf x,
\end{align}
the criterion–function reading of \textbf{Def. \ref{def:cluster_method}}, with a centroid standing in for pairwise comparison \cite{ref_10}.
The standard algorithm alternates assignment of each $\mathbf x$ to the nearest centroid with recomputation of every $\boldsymbol \mu_i$ as the mean of its members \cite{ref_5, ref_11}, as neither step increases the objective, it terminates at a local minimum only \cite{ref_2, ref_25}.

\paragraph{Hyperparameters and tuning}
\label{sec:tech:kmeans:hyperparameters}
Three: $|\mathcal C|$, fixed by the shared procedure of \textbf{Sect. \ref{sec:methodology:k}} rather than tuned per method; the initialisation and the number of restarts, which exist only because the objective is non–convex, so the best objective over several restarts is retained \cite{ref_2, ref_5}.

\paragraph{Complexity}
\label{sec:tech:kmeans:complexity}
One iteration costs $O(m |\mathcal C| n)$ time and $O(|\mathcal C| n)$ space beyond the data \cite{ref_5, ref_10}: the cheapest method surveyed, and the only one needing no pairwise dissimilarity matrix, so it remains usable at full data volume.

\paragraph{Application to \projectname{} data}
\label{sec:tech:kmeans:application}
\placeholder{Library and version, seed, restarts, and any deviation from the pipeline of Section~\ref{sec:data:preprocessing}, scaling in particular.}

\paragraph{Results}
\label{sec:tech:kmeans:results}
\placeholder{Every index of Section~\ref{sec:methodology:metrics}, the evidence behind the chosen $|\mathcal C|$, and the cluster profiles.}
% \todo{export from the benchmark notebook}

\paragraph{Limitations}
\label{sec:tech:kmeans:limitations}
Local convergence, so the partition depends on initialisation. The centroid is a mean, hence sensitive to outliers, which cannot be labelled noise. Spherical, equal–size bias splits elongated structure and merges structure of uneven density \cite{ref_2} \cite{ref_10}.
The Silhouette Coefficient of \textbf{Sect. \ref{sec:measure:silhouette}} rewards precisely the geometry K–Means produces, so that pairing flatters it \textbf{(Sect. \ref{sec:methodology:threats})}.

\paragraph{Summary}
\label{sec:tech:kmeans:summary}
The cheap, scale–sensitive baseline, linear in $m$, adequate wherever regimes are compact and comparable in size. Whether it earns more than that here is carried into \textbf{Sect. \ref{sec:comparison}}.

% ==================================================================
%  REFERENCES USED
%
%    ref_1    Singh & Singh (2024), review of clustering techniques   cited in: Overview
%    ref_2    Ezugwu et al. (2022), survey of clustering algorithms   cited in: Formulation, Hyperparameters, Limitations
%    ref_5    Xu & Tian (2015), comprehensive survey                  cited in: Overview, Formulation, Hyperparameters, Complexity
%    ref_10   Tan et al. (2019), Introduction to Data Mining, ch. 7   cited in: Overview, Assumptions, Formulation, Complexity, Limitations
%    ref_11   Han et al. (2011), Data Mining, ch. 10                  cited in: Assumptions, Formulation
%    ref_13   De Santi et al. (2025), WWTP fingerprints               cited in: Overview
%    ref_17   Guo et al. (2024), daily consumption patterns           cited in: Overview
%
%  MISSING --- the original sources are not in literature.bib and cannot
%  be cited: MacQueen (1967) and Lloyd (1982).  See the \todo in
%  `Formulation'.  Raised with the document lead.
% ==================================================================
